%! TEX root = main.tex
\begin{solution}{ex:la-7}
  (a) Yes, the map is linear. Let $u,v\in \R^2$, and $a,b\in \R$. We need to prove:
  \[
    T_1(au+bv)=aT_1(u)+bT_1(v)
  \]
  Let's check:
  \[
    \begin{gathered}
      (2(au_1+bv_1)+(au_2+bv_2),au_1+bv_1-3(au_2+bv_2))=a(2u_1+u_2,u_1-3u_2)+b(2v_1+v_2,v_1-3v_2)\\
      (2au_1+2bv_1+au_2+bv_2,au_1+bv_1-3au_2-3bv_2)=(2au_1+au_2,au_1-3au_2)+(2bv_1+bv_2,bv_1-3bv_2)\\
      (2au_1+2bv_1+au_2+bv_2,au_1+bv_1-3au_2-3bv_2)=(2au_1+2bv_1+au_2+bv_2,au_1+bv_1-3au_2-3bv_2)\\
    \end{gathered}
  \]
  Furthermore, by seeing what the image of the basis vectors $(0,1)$ and $(1,0)$ are, we can define a matrix a such that $T_1(x)=Ax$. A is defined like so:
  \[
    A=\begin{pmatrix}2&1\\1&-3\end{pmatrix}
  \]
  (b)
  No it is not linear. A linear map would satisfy $T_2(ax)=aT_2(x)$, but a counter-example with $x=(1,1)$ and $a=2$ shows:
  \[
    \begin{gathered}
      (2^2,2)\neq (2\cdot 1^2,2\cdot 1)\\
      (4,2)\neq (2,2)\\
    \end{gathered}
  \]
  (c)
  Yes it is a linear, Similar to $T_2$, except with a smaller rank. Let's prove it from the definition of linearity:
  \[
    \begin{gathered}
      T(au+bv)=aT_1(u)+bT_2(v)\\
      au_1+bv_1-2(au_2+bv_2)=a(u_1-2u_2)+b(v_1-2v_2)\\
      au_1+bv_1-2(au_2+bv_2)=au_1-2au_2)+bv_1-2bv_2)\\
      au_1+bv_1-2(au_2+bv_2)=au_1+bv_1-2au_2-2bv_2\\
      au_1+bv_1-2(au_2+bv_2)=au_1+bv_1-2(au_2+bv_2)\\
    \end{gathered}
  \]
\end{solution}
\begin{solution}{ex:la-8}
  Let's find out where $T$ puts the basis vectors $\begin{pmatrix}1\\0\end{pmatrix}$ $\begin{pmatrix}0\\1\end{pmatrix}$ using linearity. Setting $u=\begin{pmatrix}1\\1\end{pmatrix}$ and $v=\begin{pmatrix}1\\-1\end{pmatrix}$
  \[
    \begin{gathered}
      T(\frac{1}{2}u+\frac{1}{2}v)=\frac{1}{2}T(u)+\frac{1}{2}T(v)\\
      T\begin{pmatrix}1\\0\end{pmatrix}=\begin{pmatrix}\frac{1}{2}\\0\\1\end{pmatrix}+\begin{pmatrix}0\\\frac{1}{2}\\\frac{1}{2}\end{pmatrix}\\
      T\begin{pmatrix}1\\0\end{pmatrix}=\begin{pmatrix}\frac{1}{2}\\\frac{1}{2}\\\frac{3}{2}\end{pmatrix} \\\\
      T(\frac{1}{2}u-\frac{1}{2}v)=\frac{1}{2}T(u)-\frac{1}{2}T(v)\\
      T\begin{pmatrix}0\\1\end{pmatrix}=\begin{pmatrix}\frac{1}{2}\\0\\1\end{pmatrix}-\begin{pmatrix}0\\\frac{1}{2}\\\frac{1}{2}\end{pmatrix}\\
      T\begin{pmatrix}0\\1\end{pmatrix}=\begin{pmatrix}\frac{1}{2}\\-\frac{1}{2}\\\frac{1}{2}\end{pmatrix} \\
    \end{gathered}
  \]
  Now that we know where the basis vectors are in the image, we can write down a general formula for $T\begin{pmatrix}x\\y\end{pmatrix}$ and find the matrix A such that $T(x)=Ax$. The general form is:
  \[
    T\begin{pmatrix}x\\y\end{pmatrix}=(\frac{1}{2}x+\frac{1}{2}y,\frac{1}{2}x-\frac{1}{2}y,\frac{3}{2}x-\frac{1}{2}y)
  \]
  And the matrix A is:
  \[
    A=\begin{pmatrix}\frac{1}{2}&\frac{1}{2}\\\frac{1}{2}&-\frac{1}{2}\\\frac{3}{2}&\frac{1}{2}\end{pmatrix}
  \]
\end{solution}
\begin{solution}{ex:la-9}
  Let's row-reduce A:
  \[
    \begin{gathered}
      A=\begin{pmatrix}1&0&2&1\\0&1&1&2\\1&1&3&3\end{pmatrix}\\
      R_3=R_3-R_1\\
      A=\begin{pmatrix}1&0&2&1\\0&1&1&2\\0&1&1&2\end{pmatrix}\\
      R_3=R_3-R_2\\
      A=\begin{pmatrix}1&0&2&1\\0&1&1&2\\0&0&0&0\end{pmatrix}\\
    \end{gathered}
  \]
  The pivot columns are columns 1 and 2, and $Ax=0$ gives:
  \[
    \begin{gathered}
      x_1+2x_3+x_4=0\\
      x_2+x_3+2x_4=0
    \end{gathered}
  \]
  The kernel has this basis (4 unknowns with 2 equations gives 2 degrees of freedom):
  \[
    \ker A = \text{span}\left\{\begin{pmatrix}-1\\1\\1\\-1\end{pmatrix},\begin{pmatrix}2\\10\\2\\-6\end{pmatrix}\right\}
  \]
  And the image has this basis:
  \[
    \operatorname{im} A= \text{span}\left\{\begin{pmatrix}1\\0\\1\end{pmatrix},\begin{pmatrix}0\\1\\1\end{pmatrix}\right\}
  \]
  Then $\dim\ker A = 2$ and $\dim \operatorname{im}A = 2$ and $2+2=4$.
\end{solution}
\begin{solution}{ex:la-10}
  Knowing that T is injective tells us that no 2 distinct vectors have the same image under T. Now let's use the linearity of T to construct a proof by contradiction. Suppose that $T(v_1),\ldots, T(v_k)$ aren't linearly independent, i.e. there are some $v_a, v_b,v_c \in V$ and $a,b \in \R$ such that \[
    aT(v_a)+bT(v_b)=T(v_c)
  \]
  This would imply by linearity of T that
  \[
    T(av_a+bv_b)=T(v_c)
  \]
  And, since T is injective, that would mean that
  \[
    av_a+bv_b=v_c
  \]
  But since $v_1,\ldots, v_k$ is linearly independent, there's a contradiction. Thus, $T(v_1),\ldots,T(v_k)$ is linearly independent.
\end{solution}
\begin{solution}{ex:la-11}
  NOTE: reason using the dimensions, if T isn't surjective, it's rank must be smaller than dim W, so it can't be injective
\end{solution}
